
\documentclass[UTF8,zihao=-4]{ctexart}
\usepackage[a4paper,top=2.3cm,bottom=2.3cm,left=2.35cm,right=2.35cm]{geometry}
\usepackage{fontspec}
\usepackage{xeCJK}
\setmainfont{DejaVu Sans}
\setsansfont{DejaVu Sans}
\setmonofont{Noto Sans Mono}
\setCJKmainfont{Noto Sans CJK SC}
\setCJKsansfont{Noto Sans CJK SC}
\setCJKmonofont{Noto Sans CJK SC}
\usepackage{booktabs}
\usepackage{longtable}
\usepackage{array}
\usepackage{tabularx}
\usepackage{enumitem}
\usepackage{xcolor}
\usepackage{hyperref}
\usepackage{fancyhdr}
\usepackage{titlesec}
\usepackage{lastpage}
\usepackage{pifont}
\usepackage{amsmath}
\usepackage{amssymb}
\usepackage{makecell}
\usepackage{seqsplit}
\hypersetup{colorlinks=true,linkcolor=black,urlcolor=blue,citecolor=black}
\definecolor{RuleBlue}{RGB}{25,74,135}
\definecolor{LightBlue}{RGB}{236,243,252}
\definecolor{LightGray}{RGB}{247,247,247}
\definecolor{LightYellow}{RGB}{255,249,230}
\definecolor{DarkGray}{RGB}{80,80,80}
\definecolor{WarnOrange}{RGB}{172,103,0}
\definecolor{RejectRed}{RGB}{150,35,35}
\definecolor{PassGreen}{RGB}{38,118,82}
\definecolor{DiagnosticPurple}{RGB}{105,72,145}
\titleformat{\section}{\Large\bfseries\color{RuleBlue}}{\thesection}{0.8em}{}
\titleformat{\subsection}{\large\bfseries}{\thesubsection}{0.8em}{}
\titleformat{\subsubsection}{\normalsize\bfseries}{\thesubsubsection}{0.8em}{}
\setlist[itemize]{topsep=3pt,itemsep=2pt,parsep=0pt,leftmargin=1.9em}
\setlist[enumerate]{topsep=3pt,itemsep=2pt,parsep=0pt,leftmargin=2.2em}
\newcolumntype{Y}{>{\raggedright\arraybackslash}X}
\newcolumntype{L}[1]{>{\raggedright\arraybackslash}p{#1}}
\newcommand{\Pass}{\textcolor{PassGreen}{\ding{51}}}
\newcommand{\Warn}{\textcolor{WarnOrange}{\ding{115}}}
\newcommand{\Reject}{\textcolor{RejectRed}{\ding{55}}}
\newcommand{\Diag}{\textcolor{DiagnosticPurple}{\ding{72}}}
\newcommand{\Code}[1]{\texttt{#1}}
\newcommand{\MID}[1]{\texttt{\seqsplit{#1}}}
\newcommand{\Term}[1]{\textbf{#1}}
\newenvironment{notebox}[1]{\vspace{0.5em}\noindent\colorbox{LightBlue}{\parbox{0.96\linewidth}{\textbf{#1}}}\par\vspace{-0.2em}\begin{quote}\small}{\end{quote}\vspace{0.3em}}
\newenvironment{decisionbox}[1]{\vspace{0.5em}\noindent\colorbox{LightGray}{\parbox{0.96\linewidth}{\textbf{#1}}}\par\vspace{-0.2em}\begin{quote}\small}{\end{quote}\vspace{0.3em}}
\newenvironment{warningbox}[1]{\vspace{0.5em}\noindent\colorbox{LightYellow}{\parbox{0.96\linewidth}{\textbf{#1}}}\par\vspace{-0.2em}\begin{quote}\small}{\end{quote}\vspace{0.3em}}
\hypersetup{pdftitle={L3 v2 Metric Revision Update v1.1}, pdfauthor={ChatGPT for Lingji Qiwuzhi}}
\pagestyle{fancy}
\fancyhf{}
\lhead{L3 v2 Metric Revision Update v1.1}
\rhead{RDDQF / Implementation Sync}
\cfoot{第 \thepage\ 页 / 共 \pageref{LastPage} 页}
\begin{document}

\begin{titlepage}
\centering
\vspace*{2.2cm}
{\Huge\bfseries L3 v2 指标修订与实现同步说明\par}
\vspace{0.35cm}
{\Large Metric Revision and Implementation Synchronization\par}
\vspace{0.6cm}
{\LARGE v1.1\par}
\vspace{1.15cm}
\begin{tabular}{rl}
文档定位： & RDDQF v1.0 设计包的实现同步补丁文档 \\
适用系统： & 灵机启物机器人数采质检平台 / L3 v2 \\
覆盖范围： & MQ、LQ、DI、DX、Fusion、Camera Temporal Continuity \\
输入依据： & 代码集成后的指标审查、真实 wrist camera 卡顿案例、同步错位检测结果 \\
输出对象： & Agent 修改任务、后端指标实现、前端解释性 UI、报告系统 \\
核心原则： & Training Quality 与 Execution Diagnostic 严格分离；Data Integrity 具备可靠性门控作用 \\
\end{tabular}
\vfill
\begin{minipage}{0.88\linewidth}
\centering\large\bfseries
本文件不是重写 RDDQF 理论，而是把实现阶段发现的指标定义问题、时间同步问题和相机卡顿案例同步回设计文档，形成 v1.1 可执行版本。
\end{minipage}
\vfill
{\large 2026 年 6 月\par}
\end{titlepage}

\tableofcontents
\newpage

\section{v1.1 同步更新摘要}

本次同步修改主要解决六类问题：

\begin{enumerate}
\item 修正 MQ-01 将二阶差分误称为 jerk 的问题，明确 Trajectory Smoothness 应使用三阶差分或时间归一化 jerk；
\item 将 MQ-02 从 action step magnitude 修改为 action second difference continuity；
\item 将 MQ-03 从简单 sign flip P95 修改为带幅度门控的持续震荡与 hand chatter 比例；
\item 将 LQ 系列重新组织为 Coverage、Intensity、Low-value Duration 三个互补学习价值指标；
\item 将 DI 系列从单一 timestamp CV / 固定 700ms sync 阈值升级为鲁棒时间完整性与自适应同步诊断；
\item 将 DX-01 从主分数中移出，仅作为 Execution Diagnostic，并增加 lag alignment、persistent error 与 diagnostic timeline。
\end{enumerate}

同时，真实数据中出现的 wrist camera 卡顿案例说明：现有 DI-02 能捕获严重同步错位，但还需要新增 camera temporal continuity 证据，用于区分 sync offset、frame drop、stale frame reuse 和 freeze-after-jump 现象。

\section{MQ 指标修订}

\subsection{MQ-01 Trajectory Smoothness}

MQ-01 的语义目标是评价轨迹是否平滑，重点不是惩罚速度变化，而是惩罚加速度变化的突兀性。因此：

\[
jerk_t = q_t - 3q_{t-1} + 3q_{t-2} - q_{t-3}
\]

推荐主统计量为：

\[
MQ01_{raw}=P_{95}(\|jerk_t^{arm}\|_{RMS})
\]

实现要求：

\begin{itemize}
\item 旧版 \Code{np.diff(qpos, n=2)} 只能表示 acceleration-like roughness，不应再命名为 jerk；
\item MQ-01 应使用 \Code{np.diff(qpos\_arm, n=3)}，并将 timeline offset 改为 3；
\item 若未来使用非均匀 timestamp，应升级为时间归一化 jerk 或 \Code{np.gradient} 方案；
\item 原二阶差分可保留为 diagnostic feature，例如 \Code{joint\_acc\_norm}，但不再作为 smoothness 主依据。
\end{itemize}

临时评分阈值需要重新标定。初始可使用：

\[
\begin{aligned}
good &= 0.006\\
warn &= 0.018\\
bad &= 0.045
\end{aligned}
\]

\subsection{MQ-02 Motion Continuity}

MQ-02 不应评价 action 每一步走得多远，而应评价 action step 是否突然变化。新定义为：

\[
C_t=a_t-2a_{t-1}+a_{t-2}
\]

\[
MQ02_{raw}=P_{95}(\|C_t\|_{RMS})
\]

实现要求：

\begin{itemize}
\item 保留 \Code{action\_step\_norm} 给 LQ-01/LQ-02 使用；
\item 新增 \Code{action\_delta2\_norm} 给 MQ-02 使用；
\item MQ-02 timeline 使用 action second difference 异常段，不再使用一阶 action step；
\item 算法 ID 建议为 \Code{action\_second\_difference\_p95}。
\end{itemize}

初始阈值：

\[
\begin{aligned}
good &= 0.010\\
warn &= 0.030\\
bad &= 0.080
\end{aligned}
\]

\subsection{MQ-03 Motion Stability}

MQ-03 的目标是检测持续性反复修正、震荡和 hand chatter，而不是检测所有方向反转。新逻辑应加入幅度门控和 segment 持续性。

Action 有效方向反转：

\[
Rev_{t,i}=\mathbf1(|\Delta a_{t,i}|>\epsilon_i \land |\Delta a_{t-1,i}|>\epsilon_i \land sign(\Delta a_{t,i})sign(\Delta a_{t-1,i})<0)
\]

每帧反转比例：

\[
Rev_t=\frac{1}{D}\sum_i Rev_{t,i}
\]

手部 chatter 同样必须基于有效方向反转，而不是简单使用 \(|\Delta h|>2/255\)。

推荐最终定义：

\[
MQ03_{raw}=0.6R_{osc}+0.4R_{chat}
\]

其中 \(R_{osc}\) 和 \(R_{chat}\) 均为持续异常时间占比。

初始阈值：

\[
\begin{aligned}
good &= 0.05\\
warn &= 0.15\\
bad &= 0.35
\end{aligned}
\]

\section{LQ 指标修订}

\subsection{LQ-01 Effective Action Ratio}

LQ-01 应评价有效 action supervision 覆盖率。旧版只使用 episode 内部 \(P_{35}\) 容易导致整体低动态 episode 也获得较高得分。新版使用绝对阈值与动态阈值组合，并区分 arm / hand。

\[
E_t^{arm}=\mathbf1(A_t^{arm}>\max(\epsilon_{arm},P_{20}(A^{arm})))
\]

\[
E_t^{hand}=\mathbf1(A_t^{hand}>\max(\epsilon_{hand},P_{20}(A^{hand})))
\]

\[
E_t=E_t^{arm}\lor E_t^{hand}
\]

\[
LQ01=R_{eff}=\frac{1}{N}\sum_t E_t
\]

推荐 \(\epsilon_{arm}=0.003\sim0.005\) rad，\(\epsilon_{hand}=2/255\sim4/255\)。

\subsection{LQ-02 Information Density}

LQ-02 不再只是 \(P_{75}(\|\Delta a\|)+P_{75}(\|\Delta q\|)\)，而是 Coverage 与 Intensity 的组合：

\[
LQ02=0.7(R_{eff}I_a)+0.3I_q
\]

其中：

\[
I_a=P_{75}(A_t \mid E_t=1),\quad A_t=\max(A_t^{arm},A_t^{hand})
\]

\[
I_q=P_{75}(Q_t),\quad Q_t=\max(Q_t^{arm},Q_t^{hand})
\]

初始阈值：

\[
\begin{aligned}
good &= 0.030\\
warn &= 0.012\\
bad &= 0.003
\end{aligned}
\]

\subsection{LQ-03 Low-value Segment Ratio}

LQ-03 从帧级低变化比例改为连续低学习价值片段总时长占比。复用 LQ-01 的 \(E_t\)，不再重新定义 action 有效性。

\[
L_t=\mathbf1(E_t=0 \land Q_t\le \tau_Q)
\]

连续 \(L_t=1\) 的片段集合为 \(\mathcal{S}\)，每段持续时间为 \(d_k\)。

\[
LQ03=\frac{\sum_k d_k\mathbf1(d_k\ge d_{min})}{T_{episode}}
\]

推荐：

\[
d_{min}=1.0s,\quad merge\_gap=0.3s
\]

\section{DI 指标修订}

\subsection{DI-01 Timestamp Regularity}

DI-01 从 \(CV=std(dt)/mean(dt)\) 改为鲁棒时间完整性指标：

\[
DI01_{raw}=0.5J_{95}+0.3R_{gap}+0.2R_{invalid}
\]

其中：

\[
J_{95}=P_{95}\left(\frac{|dt_t-median(dt)|}{median(dt)+\epsilon}\right)
\]

\[
R_{gap}=mean(dt_t>2.0\cdot median(dt)),\quad R_{invalid}=mean(dt_t\le0)
\]

旧 timestamp CV 保留为 diagnostic evidence，不再作为主评分依据。

\subsection{DI-02 Sensor Synchronization Validity}

DI-02 从固定 700ms 阈值改为自适应同步诊断。复用 \(dt_{med}\)：

\[
\tau_{warn}=\max(0.05,2.0dt_{med})
\]

\[
\tau_{bad}=\max(0.20,6.0dt_{med})
\]

定义：

\[
R_{flag}=mean(sync\_valid=false)
\]

\[
R_{soft}=mean(sync\_diff>\tau_{warn}),\quad R_{hard}=mean(sync\_diff>\tau_{bad})
\]

\[
M_{sync}=clip(P_{95}(sync\_diff)/\tau_{bad},0,1)
\]

持续同步异常片段占比为 \(R_{seg}\)。最终：

\[
DI02_{raw}=0.30R_{flag}+0.25R_{hard}+0.20R_{soft}+0.15M_{sync}+0.10R_{seg}
\]

如果 sync metadata 缺失，不允许静默通过，应输出 \Code{sync\_metadata\_missing=true} 并至少给 warn。

\subsection{DI-03 Camera Temporal Continuity}

真实案例中，top camera 正常运动，但 wrist camera 卡在 A 点，随后突然跳到 B 点。此类问题不一定等于丢帧，更准确地说是 camera temporal discontinuity，可能来源于 frame drop、stale frame reuse、stream freeze、sync algorithm error 或 video conversion error。

建议新增 DI-03，或作为 Observation Quality 下的 OQ-01。MVP 阶段推荐归入 Data Integrity：

\[
DI03=Camera\ Temporal\ Continuity
\]

每个 camera stream 应输出：

\begin{itemize}
\item \Code{cameraId}
\item \Code{frameCount}
\item \Code{medianFps}
\item \Code{timestampGapRatio}
\item \Code{duplicateFrameRatio}
\item \Code{freezeDuringMotionRatio}
\item \Code{maxFreezeDurationSec}
\item \Code{jumpAfterFreezeCount}
\item \Code{staleFrameSegments}
\end{itemize}

推荐 raw value：

\[
DI03_{raw}=0.35R_{freeze}+0.25R_{gap}+0.20R_{duplicate}+0.20R_{jump}
\]

其中 \(R_{freeze}\) 应在机器人运动或 top camera 画面变化时检测 wrist camera 内容长期不变。若训练输入使用 wrist camera，严重 DI-03 异常应触发 reliability warning，并可参与 Data Integrity cap。

\section{DX-01 与 Fusion 修订}

\subsection{DX-01 Execution Tracking Diagnostic}

DX-01 不再进入 Training Quality Score。它用于诊断机器人执行侧是否长期偏离遥操作目标。应增加 lag alignment：

\[
k^*=\arg\min_k P_{50}(\|a_{t-k}-q_t\|),\quad k\in[0,K]
\]

推荐 \(K=5\) frames 或 150ms 至 200ms。

输出 evidence：

\[
E_{mean}=mean(E_t),\quad E_{95}=P_{95}(E_t),\quad R_{persist}=mean(E_t>\tau_{warn})
\]

Diagnostic severity：

\[
Severity_{DX01}=0.5clip(E_{95}/\tau_{bad},0,1)+0.3clip(E_{mean}/\tau_{warn},0,1)+0.2R_{persist}
\]

\[
Score_{DX01}=10(1-Severity_{DX01})
\]

并明确：

\[
DX01\notin TrainingQualityScore
\]

\subsection{Training Quality Score Fusion}

维度内部使用 mean + soft-min，避免严重短板被平均数掩盖。

\[
S_{motion}=0.8(0.35MQ01+0.35MQ02+0.30MQ03)+0.2\min(MQ01,MQ02,MQ03)
\]

\[
S_{learn}=0.8(0.25LQ01+0.35LQ02+0.40LQ03)+0.2\min(LQ01,LQ02,LQ03)
\]

\[
S_{data}=0.6(0.4DI01+0.6DI02)+0.4\min(DI01,DI02)
\]

如果纳入 DI-03，则推荐：

\[
S_{data}=0.55(0.30DI01+0.45DI02+0.25DI03)+0.45\min(DI01,DI02,DI03)
\]

主分：

\[
S_{weighted}=0.4S_{motion}+0.4S_{learn}+0.2S_{data}
\]

Data Integrity cap：

\[
S_{data}<5.0\Rightarrow S_{total}=\min(S_{weighted},6.0)
\]

\[
S_{data}<3.0\Rightarrow S_{total}=\min(S_{weighted},4.5)
\]

\section{严重同步错位的解释性输出}

严重同步错位表示同一帧质量报告中，不同数据流并不对应同一个真实物理时刻。例如：

\begin{center}
\begin{tabular}{lll}
\toprule
数据流 & 主时间轴显示 & 真实采集时间 \\
\midrule
top camera & 12.0s & 12.0s \\
wrist camera & 12.0s & 10.8s \\
qpos/action & 12.0s & 12.0s \\
\bottomrule
\end{tabular}
\end{center}

这意味着训练样本可能变为旧 wrist 画面对应当前 action，属于 observation-action temporal misalignment。

前端应增加 debug overlay：

\begin{itemize}
\item \Code{camera}: wrist\_left
\item \Code{videoTimeSec}: 12.03
\item \Code{captureTimeSec}: 10.82
\item \Code{syncDiffMs}: 1210
\item \Code{status}: severe desync
\end{itemize}

建议新增 sidecar 文件 \Code{camera\_frame\_index.json}：

\begin{quote}\small
\Code{cameraId}, \Code{frameIndex}, \Code{videoPtsSec}, \Code{cameraCaptureSec}, \Code{episodeSec}, \Code{matchedTelemetrySec}, \Code{syncDiffMs}, \Code{isStale}, \Code{isDuplicate}, \Code{sourceFrameId}
\end{quote}

这样严重同步错位不再只是一个抽象告警，而可以解释为：哪一路相机、落后多少毫秒、从第几秒开始、持续多久、是否可能是旧帧复用。

\section{Agent 修改任务总清单}

\begin{enumerate}
\item 更新 Feature Layer：新增 joint jerk、action second difference、effective action mask、robust timestamp features、adaptive sync features、camera frame temporal features；
\item 更新 Evidence Layer：输出 evidence 中必须包含 source sensor、source metric、raw values、thresholds、affected duration、timeline range；
\item 更新 Metric Layer：按本文修订 MQ-01 至 MQ-03、LQ-01 至 LQ-03、DI-01 至 DI-02、DX-01；
\item 新增 DI-03 或 OQ-01：检测 camera timestamp gap、duplicate frame、freeze during motion、jump after freeze；
\item 更新 Quality Layer：DX-01 移出主分数；DI 作为可靠性 gate；总分使用 soft-min 与 cap；
\item 更新 Timeline Engine：支持 sensor-specific timeline，例如 wrist camera severe desync、camera freeze、low-value segment、execution diagnostic segment；
\item 更新 API：增加 \Code{reliabilityWarnings}、\Code{diagnosticWarnings}、\Code{executionDiagnosticScore}、\Code{cameraFrameDiagnostics}；
\item 更新 UI：主分与执行诊断分开显示，严重同步错位需要解释受影响 sensor、offset、duration 和 training risk；
\item 保持旧 \Code{qc-context} 兼容，但新增字段应优先服务 L3 v2 页面。
\end{enumerate}

\section{最终原则}

\[
TrainingQualityScore=f(MQ,LQ,DI)
\]

\[
ExecutionDiagnosticScore=f(DX)
\]

\[
DX\notin TrainingQualityScore
\]

\[
Severe\ Temporal\ Integrity\ Issue\Rightarrow Reliability\ Warning\ and\ Score\ Cap
\]

此版本完成后，L3 v2 不再是单纯运动分数，也不是机器人控制器评分，而是面向机器人训练数据资产的可解释质量评估系统。

\end{document}
